\documentclass[12pt]{article}

\usepackage{fancyhdr}
\usepackage{enumitem}
\usepackage{amsmath}
\usepackage{amsfonts}
\usepackage{amsthm}
\usepackage{amssymb}

\renewcommand{\vec}[1]{\mathbf{#1}}

\title{Homework 3}
\author{Satvik Prasad}
\date{}

\pagestyle{fancy}
\lhead{Satvik Prasad}
\rhead{MATH3012}

\begin{document}
\begin{enumerate}
  \item We first turn the closed form of this generating function
    into a series. Using partial fractions, we get
    \begin{align*}
      f(x) = \dfrac{27}{(2x+1)(x+2)^2} = 3 \cdot
      \left[\dfrac{4}{2x+1} - \dfrac{2}{x+2}- \dfrac{3}{(x+2)^2} \right]
    \end{align*}
    Notice that
    \begin{align*}
      \dfrac{1}{x+2} = \dfrac{1/2}{1 + x/2} = \dfrac{1}{2}\sum_{k =
      0}^{\infty} (-1)^k \dfrac{x^k}{2^k}.
    \end{align*}
    Removing the $k = 0$ case from our sum and differentiating both
    sides, we get
    \begin{align*}
      - \dfrac{1}{(x+2)^2} = \sum_{k = 1}^\infty (-1)^k
      \dfrac{kx^{k-1}}{2^{k+1}} = \sum_{k = 0}^{\infty} (-1)^{k + 1}
      \dfrac{(k+1)x^{k}}{2^{k + 2}},
    \end{align*}
    meaning
    \begin{align*}
      \dfrac{1}{(x+2)^2} = \sum_{k = 0}^{\infty} (-1)^k
      \dfrac{(k+1)}{2^{k + 2}} x^k.
    \end{align*}
    Thus, we can find a closed form for $f(x)$:
    \begin{align*}
      f(x) &= 3\left[
        4\sum_{k = 0}^{\infty} (-1)^k 2^k x^k -
        \sum_{k=0}^{\infty}(-1)^k \dfrac{1}{2^k}\cdot x^k - 3\sum_{k
        = 0}^{\infty} (-1)^k \dfrac{(k+1)}{2^{k+2}}x^k
      \right] \\
      &= \sum_{k = 0}^{\infty} \left[
        3(-1)^k \cdot \left(2^{k + 2} - \dfrac{1}{2^k} - 3\cdot
        \dfrac{k + 1}{2^{k + 2}} \right)
      \right] x^k.
    \end{align*}
    Hence, the series produced by this generating function is
    \begin{align*}
      \left\{3(-1)^n \cdot \left(
          2^{n + 2} - \dfrac{1}{2^n} - \dfrac{3(n + 1)}{2^{n + 2}}
      \right)\right\}_{n \geq 0}.
    \end{align*}
    \qed

  \item The coefficient $c_p$ will be the number of solutions to the equation
    \begin{align*}
      4x_1 + 5x_2 = p,
    \end{align*}
    where $x_1, x_2 \geq 0.$ There is one way to make a sum of  $0,
    4, 8, \ldots,$ cents only the 4-cent coin, and there is one way
    to make a sum of either $0, 5, 10, \ldots,$ with only the 5-cent
    coin. The $p$-th convolution of this series will yield the number
    of ways to form a sum of $p$ with either coin. Thus, $c_p$ will
    be the coefficient of $x^p$ in the generating function
    \begin{align*}
      f(x) &= (1 + x^4 + x^8 + \ldots )(1 + x^5 + x^{10} + \ldots)  \\
      &= \dfrac{1}{1-x^4} \cdot \dfrac{1}{1-x^5}.
    \end{align*}
    Therefore, the ordinary generating function of $(c_p)_{p \geq 0}$ is
    \begin{align*}
      f(x) = \dfrac{1}{(1-x^4)(1-x^5)}.
    \end{align*}
    \qed
\end{enumerate}
\end{document}
